\section{Translation Schemes}

Consider the following generic Alloy model M:


\begin{lstlisting}
	sig A
	{
		f : A
	}
	fact R
	{
		// fact body
	}
	pred P 
	{
		// predicate body
	}

	// command C - run P
	// scope s - 2 A
	run P for 2 A

\end{lstlisting}

The core of the model M is translated as follows:

\begin{lstlisting}
	VARIABLES A, f

	R == <fact body>

	Init == 
		(A = {}
		A = {"A0"}
		\/ A = {"A0","A1"})
		/\ f \in SUBSET (A \X A)
		/\ R

	Next == UNCHANGED <<A,f>>

\end{lstlisting}



\begin{itemize}
	\item Every signature and field becomes a variable
	\item The constraints of the model are translated to appear in the Init formula
	\item The scope of the command is used to determine the allowed values for scope variables
	\item The allowed values for fields are defined in terms of the signatures that the field contains
	\item The Next formula plays no role
\end{itemize}

When TLC is run on this spec, every possible instance of the Alloy model becomes a valid initial state explored by TLC. To translate the command C, there are two possibilities:

\begin{lstlisting}

\end{lstlisting}

